\section{Metodología de simulación e implementación}
\subsection{Simulación en MATLAB}
El programa \texttt{lab7\_principal.m} compara cuatro combinaciones de
observador utilizando la misma planta, referencia y condiciones iniciales.
Se exportarán series temporales y las matrices diseñadas. Se analizarán:
la respuesta \(y[k]\), el esfuerzo \(u[k]\), la innovación, las
trayectorias estimadas, el sobrepasamiento, el tiempo de establecimiento,
el error estacionario y la activación de límites de actuación. El error de
estimación completo \(\mathbf x-\widehat{\mathbf x}\) es observable en
simulación; en el banco debe distinguirse de las mediciones físicas
disponibles.

\subsection{Arquitectura propuesta para el PSoC}
En cada evento de muestreo se adquirirá la salida de la planta mediante el
ADC, se convertirá a unidades físicas y se actualizará el estimador
seleccionado. Después se calculará la ley~\eqref{eq:control}, se limitará
el mando a la región validada y se enviará al DAC. La UART almacenará
referencia, salida, estados estimados, innovación, mando solicitado,
mando aplicado, banderas y tiempos de ejecución. La temporización se
verificará con un pin de diagnóstico y el Rigol DS1202Z-E.

El programa C inicial contiene \texttt{lab7\_estimadores.c} y
\texttt{lab7\_estimadores.h}. Es el núcleo matemático y
\emph{todavía no constituye firmware de PSoC Creator ejecutable};
resta integrar el diseño de periféricos, ISR, conversión ADC/DAC y
protecciones reales.

\subsection{Protocolo de medición}
Antes del lazo cerrado se comprobarán la continuidad de los puntos de
prueba, las masas, el ajuste de sondas \(10\times\), la salida nula,
el rango bipolar del DAC, la calibración del ADC, el periodo muestreado
y el tiempo máximo de ejecución. Se guardarán simultáneamente ondas
crudas del Rigol, registros PSoC y metadatos, con una única amplitud de
referencia física por ensayo y sin correcciones tácitas.
